\documentclass[11pt]{article}
\usepackage{amsmath,amsthm,amssymb}
\usepackage[margin=1in]{geometry}
\usepackage{booktabs}
\usepackage{array}
\usepackage{stmaryrd}
\usepackage[colorlinks=true,linkcolor=blue,citecolor=blue,urlcolor=blue]{hyperref}

\newtheorem{theorem}{Theorem}
\newtheorem{proposition}[theorem]{Proposition}
\newtheorem{corollary}[theorem]{Corollary}
\newtheorem{lemma}[theorem]{Lemma}
\theoremstyle{definition}
\newtheorem{definition}[theorem]{Definition}
\newtheorem{remark}[theorem]{Remark}
\newtheorem{example}[theorem]{Example}

\newcommand{\Cat}{\mathsf{Cat}}
\newcommand{\Set}{\mathsf{Set}}
\newcommand{\Ab}{\mathsf{Ab}}
\newcommand{\Sc}{\mathcal{S}}
\newcommand{\C}{\mathcal{C}}
\newcommand{\D}{\mathcal{D}}
\newcommand{\Oc}{\mathcal{O}}
\newcommand{\Sk}{\mathrm{Sk}}
\newcommand{\Ob}{\mathrm{Ob}}
\newcommand{\Hom}{\mathrm{Hom}}
\newcommand{\dom}{\mathrm{dom}}
\newcommand{\cod}{\mathrm{cod}}
\newcommand{\id}{\mathrm{id}}
\newcommand{\Z}{\mathbb{Z}}
\newcommand{\bowtiemap}{\bowtie}

\title{Supply-chain inventory consistency is a degree-one class\\
\large and is independent of the Zappa--Sz\'ep weld obstruction}
\author{MacBeth\thanks{Kodamai. Source proof, verification scripts, and this note live at
\texttt{github.com/scotmacbeth/work-in-progress}. Prepared for the Kodamai AI-Mathematician
grant (Applications pillar) and the Strathclyde mini-course \emph{Tangent Categories and
Containers} (13--16 October 2026).}\\
\small Kodamai}
\date{6 October 2026\\[2pt]
\small content commit \texttt{33ce5c9}}

\begin{document}
\maketitle

\begin{abstract}
A supply chain, abstracted to a small category, composes with a second chain sharing a
warehouse as a Zappa--Sz\'ep product $\Sc=\C\bowtiemap\D$, whose existence is obstructed by a
degree-two class $[\omega]\in H^2(\Sk_\Sc;\D)$. It is tempting --- and, for blockchain
re-entrancy, correct --- to identify ``inventory is inconsistent between two routes'' with
``$[\omega]\neq0$''. We show this identification fails for supply chains, in both directions.
Writing the inventory as a resource presheaf $R$ on $\Sc$, we prove that the matched-pair law
of the Zappa--Sz\'ep product is invisible to global sections of $R$ (Lemma~\ref{lem:nolambda}),
so inventory consistency is cut out entirely by the two factors and is the vanishing of a
\emph{degree-one} class $[\theta_R]\in H^1(\Sc;R)$ --- the route-reconvergence discrepancy
(Theorem~\ref{thm:degree1}). This class is logically independent of $[\omega]$: they live in
different degrees over different coefficient systems, and we exhibit explicit chains realizing
each of $([\omega]=0,[\theta_R]\neq0)$ and $([\omega]\neq0,[\theta_R]=0)$
(Theorem~\ref{thm:separation}). The conjecture was a false generalization of re-entrancy, where
the inventory \emph{is} the weld-state and no separate resource presheaf exists. The payoff is a
classification of compositional failure by cohomological degree: the degree records where the
conserved quantity lives --- on the nodes ($H^1$, route discrepancy) or in the weld ($H^2$,
handoff order).
\end{abstract}

\section{Introduction}\label{sec:intro}

Composition is where correctness is won or lost. When two procurement flows merge at a shared
warehouse, when two smart contracts interleave at a shared ledger, or when two delivery routes
reconverge on the same stock-keeping unit, the question is always the same: do the local data
glue to a single consistent global state? When they do not, the failure has an economic price
--- a double-spend, a phantom inventory, a reconciliation that never closes. A compositional
theory earns its keep when it tells you \emph{which} failures can occur and \emph{where} they
live.

The categorical model is by now standard. A supply chain, with nodes as objects and admissible
logistics operations as morphisms, is a small category; equivalently, by Ahman and Uustalu
\cite{au-dccat}, a \emph{directed container}. Composing two chains that share a set of
warehouses --- procurement $\C$ followed by logistics $\D$ --- amounts to forming a
\emph{Zappa--Sz\'ep product} $\Sc=\C\bowtiemap\D$, a strict factorization in which every
operation of $\Sc$ decomposes uniquely as a procurement step followed by a logistics step
\cite{au-distlaw,rw}. Such a factorization need not exist; its existence is governed by a
cohomological obstruction $[\omega]\in H^2(\Sk_\Sc;\D)$ in the sense of
Baues--Wirsching \cite{bw}, which vanishes exactly when the composite chain closes into the
claimed product \cite{gobstr}. In blockchain re-entrancy the same degree-two class, specialized
to $\Z/2$, \emph{is} the double-spend bit: $[\omega]=\varepsilon$, a fact we have verified in
Lean \cite{orch}.

This invites a conjecture, and it is the conjecture we test.

\begin{quote}
\textbf{(C)} Global inventory is consistent --- the resource data glue to a route-independent
stock over the whole chain --- if and only if the Zappa--Sz\'ep weld obstruction
$[\omega]\in H^2(\Sk_\Sc;\D)$ vanishes.
\end{quote}

\noindent It is a reasonable conjecture. It reads inventory inconsistency as a failure of the
chain to compose, and it is exactly what re-entrancy would lead you to expect. But it conflates
two things that a careful bookkeeping keeps apart: \emph{whether the chain factors at all} and
\emph{whether, given the operations, stock reconciles across routes}. No one had written down
the resource data as an explicit presheaf and compared its gluing obstruction to $[\omega]$.
When one does, the conjecture collapses.

\paragraph{Results.} We model inventory as a resource presheaf $R\colon\Sc^{\mathrm{op}}\to\Ab$
-- a local system of stock-groups transported along operations -- and define inventory
consistency as the existence of a global section. Three results follow.

\begin{itemize}
\item \textbf{The weld is invisible to inventory} (Lemma~\ref{lem:nolambda}). A family of local
stock values is a global section of $R$ over $\Sc=\C\bowtiemap\D$ if and only if it is
compatible with every $\C$-operation and every $\D$-operation; the matched-pair law
$\lambda$ that defines the Zappa--Sz\'ep product plays no role. Inventory consistency is the
degree-one amalgamation problem for the cover $\{\C,\D\}$ and nothing more.

\item \textbf{Inventory consistency is a degree-one class} (Theorem~\ref{thm:degree1}). Global
inventory is consistent if and only if a degree-one class $[\theta_R]\in H^1(\Sc;R)$ vanishes.
When the two factors are acyclic, Mayer--Vietoris identifies this class with the
\emph{reconvergence discrepancy}: the failure of the node-potentials determined on each factor
to agree on the shared warehouses.

\item \textbf{The two classes are independent} (Theorem~\ref{thm:separation}). Both
$([\omega]=0,\ [\theta_R]\neq0)$ and $([\omega]\neq0,\ [\theta_R]=0)$ occur. Hence neither
implication of (C) holds. The classes live in different degrees over different coefficient
systems ($R$ versus $\D$), so there is no natural map identifying them.
\end{itemize}

\paragraph{Why it looked true, and why there is no rescue.} Re-entrancy is the degenerate case
in which the conserved quantity \emph{is} the weld-state --- the composite protocol's own state,
living inside the matched pair --- so there is no separate resource presheaf and $[\omega]$ is
the whole story. Conjecture (C) silently imports this identification into supply chains, where
stock lives on the \emph{nodes} and is a genuinely separate degree-one datum
(Proposition~\ref{prop:norescue}). This also reconciles an earlier worked example
(\S\ref{sec:reconcile}): the family $W_{n,\varepsilon}$ that computes $[\omega]=\varepsilon$
measures a loop built from \emph{parallel} arrows, which closes inside a single hom-set and is
genuinely a weld phenomenon --- a correct computation of the \emph{other} class.

\paragraph{The payoff.} The honest result is more useful than (C) would have been. It
classifies compositional failure by cohomological degree, keyed to the locus of the conserved
quantity (\S\ref{sec:conclusion}): stock on the nodes gives an $H^1$ route discrepancy
(supply-chain inventory), a conserved quantity in the weld gives an $H^2$ handoff obstruction
(re-entrancy), and the factorization itself existing is a third $H^2$ condition. Two distinct
economic failure modes, cleanly separated by degree.

\paragraph{Outline.} Section~\ref{sec:prelim} recalls directed containers, the Zappa--Sz\'ep
weld obstruction, and the resource presheaf, and fixes the one modelling choice the argument
depends on. Section~\ref{sec:lemma} proves that the matched-pair law is invisible to inventory.
Section~\ref{sec:degree1} establishes the degree-one characterization and its closed form.
Section~\ref{sec:separation} proves independence. Section~\ref{sec:reconcile} explains the
provenance of (C) and reconciles the earlier model. Section~\ref{sec:conclusion} states the
degree classification and the open boundary.

\paragraph{Scope.} We fix an abstraction --- supply chain as free category on an operation
quiver, modulo imposed welds --- and analyze it honestly; we do not claim it is the unique
faithful model of any particular logistics system. Everything below is about that abstraction.

\section{Preliminaries}\label{sec:prelim}

We work with small categories and their cohomology. We recall only what the proofs use.

\subsection{Supply chains as directed containers}

A \emph{directed container} is a container equipped with the structure making its extension a
comonad; by Ahman and Uustalu \cite{au-dccat} directed containers are precisely small
categories. We use the category presentation throughout: a supply chain is a small category
$\Sc$ whose objects are nodes (warehouses, plants, retailers), whose morphisms out of a node
$a$ are the admissible logistics operations starting at $a$, and whose composition is
sequencing of operations. We write $\Sk_\Sc$ for the skeleton (one object per isomorphism
class), on which the cohomology of \S\ref{sec:prelim-weld} is computed.

\paragraph{Modelling choice (essential).} We model a supply chain as the \emph{free category on
its operation quiver}, modulo only the welds we explicitly impose. In particular two routes
$a\to d$ through distinct intermediate nodes are \emph{distinct morphisms} unless a relation
forces them equal. This is the honest object: a chain in which reconvergent routes carry
possibly different cost/stock transforms. A thin poset, by contrast, forces $rp=sq$ and makes
route discrepancy invisible by fiat; that is not a supply chain but the assertion that all
routes are already reconciled. The distinction is load-bearing: witness W1 of
Theorem~\ref{thm:separation} depends on distinct routes being distinct morphisms.

\subsection{The Zappa--Sz\'ep product and its weld obstruction}\label{sec:prelim-weld}

A \emph{Zappa--Sz\'ep product} (strict factorization) $\Sc=\C\bowtiemap\D$ consists of wide
subcategories $\C,\D\subseteq\Sc$ with $\C\cap\D=\Oc$ the discrete category on the shared object
set, such that every morphism of $\Sc$ factors \emph{uniquely} as $f=c\cdot d$ with $c\in\C$,
$d\in\D$. The product is governed by a matched-pair law, the distributive rewriting
$\lambda\colon\D\C\Rightarrow\C\D$, $d\cdot c=(c\triangleright d)\cdot(c\triangleleft d)$,
in the sense of Rosebrugh--Wood \cite{rw}; for directed containers this is the distributive law
of Ahman and Uustalu \cite{au-distlaw}.

Such a factorization need not exist. By \cite{pairwise}, $\Sc=\C\bowtiemap\D$ exists over a
chosen wide $\D$ iff a freeness condition (L) and a global-closure condition (G) hold; and by
\cite{gobstr}, under the abelian regime (H) the defect of (G) is a normalized Baues--Wirsching
2-cocycle $\omega_T$ on $\Sk_\C$ with coefficients in the vertex-group presheaf
$\D\colon\Sk_\C^{\mathrm{op}}\to\Ab$, so that
\[
  \text{(G) holds}\quad\Longleftrightarrow\quad [\omega]=0\in H^2(\Sk_\Sc;\D).
\]
Three features of $[\omega]$ matter below: it depends only on the category $\Sc$ and the chosen
right factor $\D$ (no resource data enter its definition); its coefficient system is $\D$, the
vertex groups of the invertible/weld part; and it is a \emph{degree-two} class.

\subsection{Inventory as a resource presheaf}

Inventory is data carried \emph{on top of} the category: a presheaf
$R\colon\Sc^{\mathrm{op}}\to\Ab$, where $R(a)$ is the group of stock configurations at node $a$
and $R(f)\colon R(\cod f)\to R(\dom f)$ transports stock back along the operation $f$. This is a
section of the container-logic fibration $\mathrm{Cont}(\cod)=\mathrm{Fam}(\cod^{\mathrm{op}})$
\cite{coflog}: fibrewise, a choice of stock object and its reindexing along operations.

\begin{definition}[inventory consistency]\label{def:consistency}
A \emph{global section} of $R$ is a family $(x_a\in R(a))_{a\in\Ob\Sc}$ with
$R(f)(x_{\cod f})=x_{\dom f}$ for every morphism $f$. Inventory $R$ is \emph{globally
consistent} iff it admits a global section --- a route-independent assignment of stock to every
node, compatible with every operation.
\end{definition}

We single out the case that models physical stock transport: a \emph{transport local system}
is an $R$ in which each $R(f)$ is translation by an element $\delta_f\in A$ of a fixed abelian
group $A$ (so $R(a)=A$ for all $a$, and $\delta$ is additive along composites). Then a global
section is a node-potential $\varphi\colon\Ob\Sc\to A$ with
$\delta_f=\varphi(\cod f)-\varphi(\dom f)$ for every $f$; that is, the 1-cochain $\delta$ is a
coboundary. Its obstruction is the class
\[
  [\theta_R]\ :=\ [\delta]\ \in\ H^1(\Sc;R).
\]
Contrast with $[\omega]$: $[\theta_R]$ depends on the resource presheaf $R$, not on how $\Sc$
factors; its coefficient system is $R$ (stock groups), not $\D$; and it is a \emph{degree-one}
class.

\section{The matched-pair law is invisible to inventory}\label{sec:lemma}

The first observation is that the entire Zappa--Sz\'ep structure --- the law $\lambda$ that
makes $\C\bowtiemap\D$ more than a disjoint union --- contributes nothing to the inventory
gluing problem.

\begin{lemma}[no-$\lambda$-correction]\label{lem:nolambda}
Let $\Sc=\C\bowtiemap\D$ and let $R\colon\Sc^{\mathrm{op}}\to\Set$ be any presheaf. A family
$(x_a)_{a\in\Ob\Sc}$ is a global section of $R$ over $\Sc$ if and only if it is compatible with
every $\C$-morphism and every $\D$-morphism. The matched-pair law $\lambda$ plays no role.
\end{lemma}

\begin{proof}
($\Rightarrow$) Immediate: $\C$- and $\D$-morphisms are morphisms of $\Sc$.

($\Leftarrow$) Compatibility of a family with a morphism is closed under composition. If
$R(g)(x_{\cod g})=x_{\dom g}$ and $R(h)(x_{\cod h})=x_{\dom h}$ with $\cod h=\dom g$, then by
functoriality
\[
  R(g\circ h)(x_{\cod g})=R(h)\bigl(R(g)(x_{\cod g})\bigr)
  =R(h)(x_{\dom g})=R(h)(x_{\cod h})=x_{\dom h}.
\]
Since every morphism of $\Sc$ is a word in $\C$- and $\D$-arrows --- indeed $f=c\cdot d$ by the
factorization --- a family compatible with all $\C$- and $\D$-generators is compatible with all
morphisms. The rewriting $\lambda$ converts $\D\C$-words to $\C\D$-words but never changes the
\emph{morphism} it names, hence never changes the compatibility constraint.
\end{proof}

\begin{remark}
The point is not that $\lambda$ is trivial --- it is generally nontrivial, and it is exactly
what $[\omega]$ measures --- but that it is inert \emph{on sections}. The weld rearranges how a
morphism is presented; a presheaf only sees the morphism. This is why no amount of weld
structure can inject content into the inventory gluing problem, as long as $\Sc$ is a
well-defined category at all.
\end{remark}

\begin{remark}[verification]
Lemma~\ref{lem:nolambda} was brute-checked on a genuinely nontrivial Zappa--Sz\'ep product ---
the flip matched pair on two objects, $t_1\cdot c=c\cdot t_0$ with $\lambda\neq\id$ --- over all
functorial $R$ valued in $\Z/3$: the global sections equalled the $\{\C,\D\}$-generator-%
compatible families, in every case (\texttt{scratch/zs\_presheaf\_sections.py}).
\end{remark}

\noindent\textbf{Consequence.} The set of consistent inventories is cut out entirely by the
$\C$- and $\D$-constraints. Inventory consistency is exactly the degree-one amalgamation problem
for the cover $\{\C,\D\}$ --- which is what the next section computes.

\section{Inventory consistency is a degree-one class}\label{sec:degree1}

\begin{theorem}\label{thm:degree1}
Let $\Sc=\C\bowtiemap\D$ (so the decomposition exists, $[\omega]=0$) and let $R$ be a transport
local system of abelian stock-groups. Then global inventory is consistent if and only if
$[\theta_R]=0\in H^1(\Sc;R)$. If moreover the factors are acyclic ---
$H^1(\C;R)=H^1(\D;R)=0$, the ``tree-like'' procurement and delivery case --- then
Mayer--Vietoris for the cover $\{\C,\D\}$ gives the closed form
\[
  H^1(\Sc;R)\ \cong\ \operatorname{coker}\!\bigl[\,H^0(\C;R)\oplus H^0(\D;R)\ \to\ H^0(\Oc;R)\,\bigr],
\]
so the obstruction is exactly the \emph{reconvergence discrepancy}: the failure of the
node-potentials determined on each factor to agree on the shared nodes.
\end{theorem}

\begin{proof}
By Lemma~\ref{lem:nolambda} a global section is a family satisfying the $\C$- and
$\D$-constraints. For a transport local system this says precisely that there is a
node-potential $\varphi$ with $\delta=d^0\varphi$, i.e.\ the 1-cochain $\delta$ is a coboundary,
i.e.\ $[\theta_R]=[\delta]=0$ in $H^1(\Sc;R)$. This equivalence uses only Lemma~\ref{lem:nolambda}
and is unconditional on the factors.

For the closed form: the cover $\{\C,\D\}$ with discrete overlap $\Oc$ presents $\Sc$ as a
pushout of categories glued along $\Oc$, then welded by $\lambda$, which by
Lemma~\ref{lem:nolambda} is inert on cohomology of coefficient presheaves. The Mayer--Vietoris
sequence of this cover reads
\[
  0\to H^0(\Sc)\to H^0(\C)\oplus H^0(\D)\to H^0(\Oc)\to H^1(\Sc)\to H^1(\C)\oplus H^1(\D)\to\cdots
\]
(coefficients $R$ throughout). The overlap $\Oc$ is discrete, so $H^{\geq1}(\Oc;R)=0$ and the
two-piece cover with discrete overlap has no higher \v{C}ech terms; acyclic factors kill
$H^1(\C)\oplus H^1(\D)$. What remains is
$H^1(\Sc;R)\cong\operatorname{coker}[H^0(\C)\oplus H^0(\D)\to H^0(\Oc)]$.
\end{proof}

\begin{remark}
Only the \emph{closed form} needs acyclic factors. The equivalence ``consistent $\iff
[\theta_R]=0$'' holds for any factors, since it rests only on Lemma~\ref{lem:nolambda}. For
factors with their own loops, $H^1(\Sc;R)$ acquires the $H^1(\C)\oplus H^1(\D)$ summand and the
decomposition changes; the characterization does not.
\end{remark}

\begin{example}[the reconvergent diamond]\label{ex:diamond}
Take the free category on the diamond quiver $a\rightrightarrows\{b,c\}\rightrightarrows d$:
routes $a\to b\to d$ (arrows $p,r$) and $a\to c\to d$ (arrows $q,s$), kept distinct. The
classifying space is the diamond graph, with $\beta_1=4-4+1=1$, so $H^1(\Sc;A)=A$ for any
coefficient group $A$; the single class is the loop-sum
$\delta_p+\delta_r-\delta_q-\delta_s$. With transport deltas $(\delta_p,\delta_r,\delta_q,\delta_s)
=(2,1,1,3)$ the holonomy is $2+1-1-3=-1\neq0$: the stock delivered at $d$ differs by route, so
inventory is inconsistent, and the class is the generator of $H^1$. A coboundary
$\delta=d^0\varphi$ is consistent (the diamond case is checked in
\texttt{scratch/supply\_chain\_obstruction.py}). The failure is \emph{degree one}, unambiguously.
Over $\Z/n$ the loop-sum lives in $\Z/n$ and consistency is the vanishing of its residue; for a
directed $k$-cycle with uniform step this is the familiar $H^1(C_k;\Z/n)=\Z/\gcd(k,n)$.
\end{example}

\section{The two classes are independent}\label{sec:separation}

\begin{theorem}[separation]\label{thm:separation}
The classes $[\theta_R]\in H^1(\Sc;R)$ and $[\omega]\in H^2(\Sk_\Sc;\D)$ are logically
independent. Both of the following occur:
\begin{enumerate}
\item[\textup{(W1)}] $[\omega]=0$ with $[\theta_R]\neq0$: the chain decomposes, yet inventory is
inconsistent.
\item[\textup{(W2)}] $[\omega]\neq0$ with $[\theta_R]=0$: the chain does not decompose, yet
inventory is trivially consistent.
\end{enumerate}
Hence neither implication of (C) holds. Moreover the two classes live in different degrees and
over different coefficient systems ($R$ versus $\D$).
\end{theorem}

\begin{proof}
\emph{Witness W1 (breaks ``$\Leftarrow$'').} Let $\Sc$ be the free category on the diamond
quiver of Example~\ref{ex:diamond}. The quiver is directed-acyclic, so $\Sc$ has no
non-identity invertible morphisms; the weld coefficient system is the trivial group, $\D=0$,
whence $H^2(\Sk_\Sc;\D)=H^2(\Sk_\Sc;0)=0$ and $[\omega]=0$ is forced. (A free category is a
strict factorization with nothing to obstruct.) Yet $\beta_1=1$, so $H^1(\Sc;R)=R$; choosing the
transport system $(2,1,1,3)$ of Example~\ref{ex:diamond} gives $[\theta_R]\neq0$ --- inventory
inconsistent.

\emph{Witness W2 (breaks ``$\Rightarrow$'').} Let $\Sc$ be the rigid-twist category of
\cite{gobstr}: a branch $a\to x\rightrightarrows y$ with $\Hom(y,y)=\Z/2$ and the twist
$s_2\circ p=(s\circ p)\cdot g$, for which $[\omega]$ is the generator of $H^2(\Sk_\Sc;\Z/2)=\Z/2$,
so $[\omega]\neq0$. Equip it with any coboundary inventory $\delta=d^0\varphi$ --- e.g.\ the zero
transport system. Such an inventory is globally consistent by construction, $[\theta_R]=0$,
while $[\omega]\neq0$.

Finally, in W1 the weld coefficient $\D$ is forced to $0$ while the resource coefficient $R$ is
nonzero; the classes cannot even be compared by a natural map, let alone identified. Both
witnesses are machine-checked in \texttt{scratch/separation\_witnesses.py}.
\end{proof}

\begin{remark}[reading the two classes]
$[\omega]$ answers \emph{``does the supply chain even factor as procurement $\bowtiemap$
logistics, as the model asserts?''} --- a structural precondition on the category. $[\theta_R]$
answers \emph{``given the operations, does stock glue to a route-independent inventory?''} ---
the actual consistency question, about data on top of the category. Conjecture (C) conflated a
precondition in degree two with the consistency class in degree one.
\end{remark}

\section{Where (C) came from, and why nothing rescues it}\label{sec:reconcile}

The conjecture is not arbitrary; it is a false generalization of blockchain re-entrancy, where
$[\omega]=\varepsilon\in H^2\cong\Z/2$ is the double-spend bit \cite{orch}. It is worth stating
precisely why the template does not transfer.

\begin{proposition}[no rescue regime]\label{prop:norescue}
There is no non-degenerate supply chain in which inventory consistency of a resource presheaf
$R$ literally equals $[\omega]=0$. The appearance of equivalence in re-entrancy is not a special
case of (C); it is a different problem, in which there is no separate resource presheaf at all.
\end{proposition}

\begin{proof}
Inventory consistency is $[\theta_R]=0$ in degree one with coefficients $R$
(Theorem~\ref{thm:degree1}); $[\omega]=0$ is a degree-two condition with coefficients $\D$. An
equivalence would require a natural identification $H^1(\Sc;R)\cong H^2(\Sk_\Sc;\D)$ of the two
classes. No such identification exists in general: there is no suspension or dimension-shift
relating a degree-one presheaf class to a degree-two weld class with different coefficients. For
the free loop $\Sc=B\Z$ one has $H^1(B\Z;A)=A$ while $H^2(B\Z;A)=0$ --- the degrees simply do not
match. And Theorem~\ref{thm:separation} already shows the two classes vary independently.
\end{proof}

\paragraph{What re-entrancy actually is.} In re-entrancy the conserved quantity \emph{is the
weld-state} --- the composite protocol's own state, living inside the matched pair, not on a
separate presheaf of nodes. There $[\omega]$ is the whole story because there is nothing else:
inventory equals weld-state by construction, so its consistency obstruction \emph{is} the weld
$H^2$. Conjecture (C) imported this identification --- ``inventory = weld-state'' --- into supply
chains, where it is false: supply-chain stock lives on the nodes, transported by operations, and
is a genuinely separate degree-one datum.

\subsection*{Reconciliation with the $W_{n,\varepsilon}$ model}\label{sub:reconcile}

An earlier worked example \cite{supplyzs} models ``inventory inconsistency between two routes''
by the warehouse family $W_{n,\varepsilon}$: objects $Wh,Pr,De$, with $s\cdot p=q$ and
$s_2\cdot p=q\cdot\tau^\varepsilon$, computing $[\omega]=\varepsilon\in H^2(\Sk_\C;\Z/n)$. This
computation is \emph{correct for the shape it uses} --- but the shape decides which of the two
consistency notions it captures, and it is not the one the phrase ``two routes'' most naturally
names.

Look at the geometry. In $W_{n,\varepsilon}$ the ``two routes'' are the \emph{parallel} arrows
$s,s_2$ sharing the same source and target (a rigid twist $a\to x\rightrightarrows y$); they
differ by a vertex-group element $\tau^\varepsilon$. A loop built from parallel arrows closes
inside a single hom-set, so its class lives in the \emph{vertex group}, hence in
$H^2(\Sk;\D)$. That is genuinely a weld phenomenon --- whether the lot-cursor algebra closes
into a distributive law. Call it \emph{provenance consistency}; it is a legitimate degree-two
question. Contrast the honest ``two routes to a delivered good'': $a\to b\to d$ and $a\to c\to d$
through \emph{distinct} intermediate nodes (the diamond, Example~\ref{ex:diamond}). Here the loop
runs through the nerve of $\Sc$, not a single hom-set; its class lives in $H^1(\Sc;R)$. This is
node-stock consistency --- degree one.

So the dividing line is clean, decided entirely by \emph{where the reconvergence loop lives}:
\begin{center}
\begin{tabular}{@{}lll@{}}
\toprule
Reconvergence loop & Class & Reads \\
\midrule
parallel arrows (vertex-group loop) & $H^2(\Sk;\D)$ & weld / provenance \\
distinct-node reconvergence (nerve loop) & $H^1(\Sc;R)$ & node stock \\
\bottomrule
\end{tabular}
\end{center}
The $W_{n,\varepsilon}$ computation therefore produced the weld-provenance class and labelled it
``inventory inconsistency between two routes.'' Under the natural reading of that phrase ---
distinct routes, reconvergent stock --- the obstruction is $H^1$, and
Theorem~\ref{thm:separation} shows it is independent of the $H^2$ the example computed. The
example is not wrong; it is a correct computation of the \emph{other} class, and should be
relabelled provenance consistency.

\section{The degree classification}\label{sec:conclusion}

Conjecture (C) is false: inventory consistency is not the vanishing of the Zappa--Sz\'ep weld
class. It is the vanishing of a degree-one class $[\theta_R]\in H^1(\Sc;R)$ on the resource
presheaf --- the route-reconvergence discrepancy --- which is independent of the degree-two weld
class $[\omega]\in H^2(\Sk_\Sc;\D)$ in degree, in coefficients, and in truth value.

The honest result organizes compositional failure by cohomological degree, keyed to the locus of
the conserved quantity.

\begin{center}
\begin{tabular}{@{}lll@{}}
\toprule
Conserved quantity lives\ldots & Failure class & Economic instance \\
\midrule
on the \textbf{nodes} (transported by operations) & $[\theta_R]\in H^1$ --- route discrepancy & supply-chain inventory \\
in the \textbf{weld} (the composite's own state) & $[\omega]\in H^2$ --- handoff order & blockchain re-entrancy \\
in the \textbf{factorization existing at all} & $[\omega]\in H^2$ --- transversal closure & does the chain decompose? \\
\bottomrule
\end{tabular}
\end{center}

\noindent Supply chains are therefore not a second instance of ``failure $=$ nonzero $H^2$.''
They are the first clean instance of ``failure $=$ nonzero $H^1$'': a route/reconvergence
discrepancy, a distinct and arguably more common economic failure mode. The degree of the
obstruction tells you where to look for the lost quantity.

\paragraph{Open boundary.} Three honest limits.
\begin{enumerate}
\item \emph{Torsor-valued inventory.} Lemma~\ref{lem:nolambda} and Theorem~\ref{thm:degree1} are
for inventory valued in a sheaf of \emph{groups} (stock levels). If inventory is valued in
\emph{torsors} --- stock defined only up to a gauge or allocation choice --- gluing becomes a
gerbe problem and the obstruction rises to $H^2$. A genuine degree-two inventory obstruction
then exists; but it is a \emph{higher-gauge} class with coefficients in $R$, still not equal to
the weld $[\omega]$ (different coefficients). Characterizing it, and its relation to $[\omega]$,
is the natural follow-up --- very likely another ``rhyme, not identity.''
\item \emph{Non-acyclic factors.} The closed form in Theorem~\ref{thm:degree1} assumes
$H^1(\C)=H^1(\D)=0$. For factors with their own loops the equivalence
``consistent $\iff[\theta_R]=0$'' still holds (Lemma~\ref{lem:nolambda} is unconditional); only
the explicit decomposition acquires the $H^1(\C)\oplus H^1(\D)$ summand.
\item \emph{The general comparison map.} Proposition~\ref{prop:norescue} rules out a natural
identification via explicit degree mismatch and the separation witnesses. The general question
--- when a connecting map in the descent spectral sequence of the $\{\C,\D\}$ cover relates
node-$H^1$ to weld-$H^2$ --- is sketched, not developed, and is not needed for the main theorems.
\end{enumerate}

\paragraph{Fidelity.} Throughout, ``a supply chain is this category'' is a modelling
commitment, not a claim that the abstraction is uniquely faithful to any real logistics system;
object-level fidelity remains open. The mathematics above concerns the abstraction we fixed in
\S\ref{sec:prelim}.

\paragraph{For the grant.} The Applications/Impact story is sharpened, not weakened: two
compositional failure modes, cleanly separated by cohomological degree, each with an economic
witness and a machine-checkable obstruction. Correctness of composition has a price, and the
degree of the obstruction tells you which account it is charged to.

\begin{thebibliography}{9}

\bibitem{au-dccat}
D.~Ahman, T.~Uustalu,
\emph{Directed containers as categories},
Proc.\ MSFP 2016, EPTCS 207, 2016, 89--98; arXiv:1604.01187.

\bibitem{au-distlaw}
D.~Ahman, T.~Uustalu,
\emph{Distributive laws of directed containers},
Progress in Informatics \textbf{10} (2013), 3--18, DOI 10.2201/NIIPI.2013.10.2.

\bibitem{bw}
H.-J.~Baues, G.~Wirsching,
\emph{Cohomology of small categories},
J.\ Pure Appl.\ Algebra \textbf{38} (1985), 187--211.

\bibitem{rw}
R.~Rosebrugh, R.~J.~Wood,
\emph{Distributive laws and factorization},
J.\ Pure Appl.\ Algebra \textbf{175} (2002), 327--353.

\bibitem{weibel}
C.~A.~Weibel,
\emph{An Introduction to Homological Algebra},
Cambridge Studies in Advanced Mathematics \textbf{38}, Cambridge Univ.\ Press, 1994.

\bibitem{pairwise}
MacBeth,
\emph{The pairwise Zappa--Sz\'ep criterion}
(Kodamai internal note, 10 June 2026; proved).

\bibitem{gobstr}
MacBeth,
\emph{The global-closure obstruction $(G)$ is a cohomology class}
(Kodamai internal note, 11 June 2026; proved).

\bibitem{orch}
MacBeth,
\emph{The re-entrancy obstruction is the token-mutation bit: an analytic proof for orchestration
Zappa--Sz\'ep products}
(Kodamai internal note, 20 July 2026; proved, $[\omega]=\varepsilon$ core Lean-verified).

\bibitem{supplyzs}
MacBeth,
\emph{A supply chain is a directed container, and composing two is a Zappa--Sz\'ep product: a
computed worked example}
(Kodamai internal note, 23 July 2026; computed).

\bibitem{coflog}
MacBeth,
\emph{The logic of containers: $\mathrm{Cont}(\cod)$ is the bifibration
$\mathrm{Fam}(\cod^{\mathrm{op}})$}
(Kodamai internal note, 2026; proved).

\end{thebibliography}

\end{document}
